\documentclass[11pt,a4paper]{article}

% ---------- Codifica e lingua ----------
\usepackage[T1]{fontenc}
\usepackage[utf8]{inputenc}
\usepackage[main=italian,provide=*]{babel}

% ---------- Matematica ----------
\usepackage{amsmath,amssymb}
\usepackage{mathtools}
\usepackage{nicematrix}

% ---------- Impaginazione ----------
\usepackage[a4paper,margin=2.7cm]{geometry}
\usepackage{caption}
\usepackage{booktabs}
\usepackage{enumitem}
\usepackage{placeins}
\usepackage[colorlinks=true,linkcolor=black,citecolor=black,urlcolor=blue!55!black]{hyperref}

\newcommand{\ket}[1]{\lvert #1 \rangle}
\newcommand{\bra}[1]{\langle #1 \rvert}
\newcommand{\Tr}{\operatorname{Tr}}

\title{Parte 2, Passo 1 --- costruzione del modello di rumore\\
\large Dall'operatore densità al \texttt{NoiseModel} di Qiskit}
\author{Samuele \\ Relatore: Prof.\ Alessandro Chiesa}
\date{}

\begin{document}
\maketitle
\tableofcontents

\bigskip
\noindent Questo documento deriva, passo per passo, tutto ciò che serve per costruire
il modello di rumore del dimero: dal motivo per cui uno stato puro non basta più, fino
alla formula esatta che converte l'errore medio di gate riportato nelle calibrazioni
pubblicate nel parametro che Qiskit vuole in ingresso. Segue la stessa notazione delle
slide del corso (\texttt{7Physical\_\allowbreak implementation.pdf}), citate dove serve.

\section{Perché serve l'operatore densità}

In un sistema chiuso lo stato è un vettore $\ket{\psi}$ e l'evoluzione è unitaria:
$\ket{\psi(t)}=U(t)\ket{\psi(0)}$. Questo smette di bastare quando il sistema $S$ (i
nostri qubit) interagisce con un ambiente $E$ che non osserviamo (il resto del
dispositivo, i suoi gradi di libertà termici, l'apparato di misura). Lo stato
composito $SE$ resta puro, ma non lo è più lo stato del solo $S$: per descriverlo si
deve tracciare via l'ambiente,
\begin{equation}
\rho^S = \Tr_E\!\big[\rho^{SE}\big],
\end{equation}
e $\rho^S$ così ottenuto è in generale uno stato \emph{misto}. Questo è esattamente il
meccanismo della decoerenza discusso a lezione: se due stati del sistema si
correlano con stati ortogonali dell'ambiente, l'informazione di fase fra di essi va
persa localmente, anche se il sistema composito resta in uno stato puro.

Un esempio minimo, utile a fissare le idee: sia $\ket{\Psi}=\alpha\ket{0_A0_E}+
\beta\ket{1_A0_B}$ (spin $A$ correlato con un secondo sistema $B$ tramite
l'entanglement). Tracciando via $B$,
\begin{equation}
\rho^A = \Tr_B\ket{\Psi}\bra{\Psi} = |\alpha|^2\ket{0_A}\bra{0_A} +
|\beta|^2\ket{1_A}\bra{1_A},
\end{equation}
che è diagonale: tutta l'informazione sulla fase relativa fra $\alpha$ e $\beta$ è
sparita da $\rho^A$, pur essendo lo stato globale $\ket{\Psi}$ perfettamente puro e
coerente. Da qui in avanti lavoriamo quindi con $\rho$, non più con $\ket{\psi}$; in
Qiskit questo corrisponde a passare da \texttt{Statevector} a
\texttt{DensityMatrix}/\texttt{AerSimulator(method="density\_matrix")}.

\section{Operazioni quantistiche e operatori di Kraus}
\label{sec:kraus}

Un'evoluzione non unitaria di $\rho$ (dovuta all'accoppiamento con $E$) si scrive in
generale come un'operazione quantistica
\begin{equation}
\mathcal E(\rho) = \sum_k E_k\,\rho\,E_k^\dagger, \qquad
\sum_k E_k^\dagger E_k = \mathbb{I},
\label{eq:kraus}
\end{equation}
dove gli $E_k$ sono gli \emph{operatori di Kraus} del canale. La condizione di
completezza $\sum_k E_k^\dagger E_k=\mathbb{I}$ garantisce $\Tr[\mathcal E(\rho)]=1$
per ogni $\rho$ normalizzato, cioè che $\mathcal E(\rho)$ resti un operatore densità
legittimo. Questa è la forma generale da cui discendono tutti i canali di rumore usati
a lezione (defasamento, rilassamento, depolarizzazione): cambia solo la scelta degli
$E_k$.

\emph{Nota terminologica}: nella tesi si useranno sempre $\varepsilon$ per l'errore
medio di gate (quantità sperimentale, da randomized benchmarking) e $\lambda$ per il
parametro del canale depolarizzante di Qiskit (quantità che entra nella
Sez.~\eqref{eq:kraus}) --- non sono la stessa cosa, si veda la Sez.~\ref{sec:eps-lambda}.

\section{Il canale depolarizzante}
\label{sec:depol}

\subsection{Un qubit}

Il canale depolarizzante a un qubit sostituisce lo stato, con probabilità $\lambda$,
con lo stato massimamente misto $\mathbb{I}/2$, lasciandolo inalterato con probabilità
$1-\lambda$:
\begin{equation}
\mathcal E(\rho) = (1-\lambda)\,\rho + \lambda\,\frac{\mathbb{I}}{2}.
\end{equation}
Usando $\mathbb{I} = (\rho+X\rho X+Y\rho Y+Z\rho Z)/4$ per ogni $\rho$ normalizzato
(identità standard per un singolo qubit, si veda Nielsen \& Chuang, citata anche nella
slide del corso), si riscrive come combinazione dello stato originale e dei tre stati
ruotati da un errore di Pauli:
\begin{align}
\mathcal E(\rho) &= (1-\lambda)\rho + \lambda\,\frac{\rho+X\rho X+Y\rho Y+Z\rho Z}{4}
= \Big(1-\tfrac{3\lambda}{4}\Big)\rho +
\frac{\lambda}{4}\big(X\rho X+Y\rho Y+Z\rho Z\big).
\end{align}
In forma di Kraus (Eq.~\eqref{eq:kraus}) questo è esattamente un canale a 4 operatori,
\begin{equation}
E_0=\sqrt{1-\tfrac{3\lambda}{4}}\,\mathbb{I},\quad
E_1=\sqrt{\tfrac{\lambda}{4}}\,X,\quad
E_2=\sqrt{\tfrac{\lambda}{4}}\,Y,\quad
E_3=\sqrt{\tfrac{\lambda}{4}}\,Z,
\label{eq:kraus-depol-1q}
\end{equation}
e si verifica direttamente che $\sum_k E_k^\dagger E_k=\mathbb{I}$ (verificato
numericamente per diversi valori di $\lambda$, residuo a precisione macchina).

\subsection{$n$ qubit}

La stessa costruzione si generalizza a $n$ qubit sostituendo lo stato mescolato con
l'identità sull'intero spazio a $d=2^n$ dimensioni:
\begin{equation}
\mathcal E(\rho) = (1-\lambda)\,\rho + \lambda\,\frac{\mathbb{I}}{2^n}.
\label{eq:depol-nqubit}
\end{equation}
È questo il canale che useremo per gli errori di gate: $n=1$ per i gate a un qubit
($R_z$, $\sqrt X$, $X$), $n=2$ per i CNOT.

\section{Errore medio di gate $\varepsilon$ e parametro $\lambda$}
\label{sec:eps-lambda}

Le calibrazioni pubblicate dai produttori di hardware non riportano $\lambda$: riportano
un \emph{errore medio di gate} $\varepsilon$, misurato con tecniche di randomized
benchmarking. Serve quindi la relazione fra le due quantità, che deriviamo qui invece
di darla per buona.

\subsection{Fedeltà media di un canale}

La fedeltà media di un'operazione quantistica $\mathcal E$ rispetto all'identità è
definita come la media, su tutti gli stati puri $\ket\psi$ dello spazio a $d$
dimensioni (con la misura di Haar), della sovrapposizione fra lo stato in ingresso e
quello in uscita:
\begin{equation}
F_\text{avg}(\mathcal E) = \int d\psi\;\bra\psi\,\mathcal E(\ket\psi\bra\psi)\,\ket\psi .
\end{equation}
L'errore medio di gate è per definizione il suo complemento, $\varepsilon \equiv
1-F_\text{avg}$.

\subsection{Calcolo per il canale depolarizzante}

Per il canale depolarizzante a $n$ qubit, Eq.~\eqref{eq:depol-nqubit}, il calcolo si fa
esplicitamente. Per ogni stato puro $\ket\psi$ normalizzato,
\begin{equation}
\bra\psi\,\mathcal E(\ket\psi\bra\psi)\,\ket\psi =
(1-\lambda)\underbrace{\bra\psi\psi\rangle\langle\psi\ket\psi}_{=1} +
\lambda\,\bra\psi\frac{\mathbb{I}}{d}\ket\psi =
(1-\lambda) + \frac{\lambda}{d},
\end{equation}
dove si è usato $\bra\psi\mathbb{I}\ket\psi=1$. Il risultato non dipende da
$\ket\psi$: l'integrale su $\psi$ è quindi banale, e
\begin{equation}
F_\text{avg} = (1-\lambda) + \frac{\lambda}{d} = 1-\lambda\Big(1-\frac1d\Big)
= 1-\lambda\,\frac{d-1}{d}.
\end{equation}
Quindi
\begin{equation}
\boxed{\;\varepsilon = 1-F_\text{avg} = \lambda\,\frac{d-1}{d}\;}, \qquad d=2^n.
\label{eq:eps-lambda}
\end{equation}

\subsection{Conversione per 1 e 2 qubit}

Invertendo l'Eq.~\eqref{eq:eps-lambda} per i due casi che useremo:
\begin{align}
n=1,\ d=2:&\quad \varepsilon = \lambda\cdot\frac12 \;\Longrightarrow\; \lambda = 2\varepsilon, \\[4pt]
n=2,\ d=4:&\quad \varepsilon = \lambda\cdot\frac34 \;\Longrightarrow\; \lambda = \frac43\,\varepsilon.
\end{align}
Sono le due conversioni confermate dal relatore e già verificate numericamente
(\texttt{depolarizing\_\allowbreak error} di Qiskit Aer riproduce $\varepsilon$ a
precisione macchina se alimentato con $\lambda=2\varepsilon$ o
$\lambda=\tfrac43\varepsilon$,
controllato contro \texttt{qiskit.quantum\_info.\allowbreak average\_gate\_fidelity},
non solo contro la formula).

\section{Errore di lettura (readout)}

Il secondo canale confermato riguarda la misura, non l'evoluzione: con probabilità $p$
il risultato letto è quello sbagliato. Seguendo la notazione della slide del corso, i
proiettori ideali $P_0=\ket0\bra0$, $P_1=\ket1\bra1$ vengono sostituiti da versioni
``deformate'',
\begin{equation}
\pi_0 = (1-p)\,P_0 + p\,P_1, \qquad \pi_1 = (1-p)\,P_1 + p\,P_0,
\end{equation}
che nel caso simmetrico (stessa probabilità di errore per $0\to1$ e $1\to0$, la scelta
richiesta dal relatore come primo passo) corrispondono alla matrice di confusione
\begin{equation}
\begin{NiceArray}{cc}[first-row,first-col]
 & \ket0\ \text{letto} & \ket1\ \text{letto} \\
\ket0\ \text{vero} & 1-p & p \\
\ket1\ \text{vero} & p & 1-p
\end{NiceArray}
\end{equation}
che è l'oggetto che si passa direttamente a \texttt{ReadoutError} in Qiskit Aer.

\section{Cosa non entra in questo modello}

Le slide del corso presentano anche il canale di defasamento (\emph{phase damping},
$T_2$) e quello di rilassamento (\emph{amplitude damping}, $T_1$), ciascuno costruito
allo stesso modo a partire da un'interazione esplicita con un ambiente a due livelli.
Per decisione del relatore questi due canali sono esclusi dallo scope della tesi: il
modello di rumore userà solo i due canali derivati sopra (gate, Sez.~\ref{sec:depol};
readout, sezione precedente). Li si menziona qui solo per completezza, non perché
verranno implementati.

\section{Verso il codice: la struttura del \texttt{NoiseModel}}

Con le Eq.~\eqref{eq:eps-lambda} e la matrice di confusione della sezione precedente,
la costruzione del modello in Qiskit è diretta: si converte l'$\varepsilon_{1q}$ del
chip scelto in $\lambda_{1q}=2\varepsilon_{1q}$ e lo si applica ai gate a un qubit
($R_z$, $\sqrt X$, $X$); si converte l'$\varepsilon_{2q}$ in
$\lambda_{2q}=\tfrac43\varepsilon_{2q}$ e lo si applica al CNOT; si costruisce la
matrice di confusione dal $p_\text{readout}$ del chip e la si applica alla misura
finale. Il limite di rumore nullo ($\lambda_{1q}=\lambda_{2q}=p_\text{readout}=0$) deve
riprodurre esattamente i risultati della Parte 1 --- è il primo controllo da fare prima
di usare il modello su un circuito completo, e va rifatto ogni volta che si cambia
qualcosa a monte nella pipeline.

\section{Prossimo passo}

Applicare questo \texttt{NoiseModel}, con valori di calibrazione reali e citabili, alla
stima VQE del ground state con termine DM (Passo 2 della pipeline).

\end{document}
